\chapter{Top-Level Datapath}

\section{Introduction}

The datapath forms the computational backbone of the Sargantana processor. It interconnects
every major pipeline stage, coordinates the movement of instructions and operands, and
interfaces with all execution units. While the control logic determines \textit{what} operation must be
performed, the datapath determines \textit{where} information flows and \textit{how} the processor state is
updated.

At the highest level, the datapath receives instructions from the instruction fetch subsystem,
processes them through decoding, register renaming, operand reading, execution, and finally
commits the architectural state once execution completes successfully.

The complete organization of the datapath is illustrated in Figure~\ref{fig:datapath-overview}.

\begin{figure}[H]
  \centering
  \includegraphics[width=0.85\textwidth]{ch03_datapath_overview.png}
  \caption{Top-level organization of the Sargantana datapath}
  \label{fig:datapath-overview}
\end{figure}

\section{Top-Level Interfaces}

The top-level datapath module serves as the central connection point between the frontend,
backend, execution units, memory subsystem, and commit logic.

Its primary interfaces include:

\subsection*{Instruction Interface}
Receives decoded instructions originating from the frontend pipeline.
Connected modules include:
\begin{itemize}
  \item Decoder
  \item Instruction Queue
\end{itemize}

\subsection*{Register Interface}
Communicates with the physical register file to obtain source operands and write execution
results.
Connected modules include:
\begin{itemize}
  \item Register File
  \item Rename Logic
  \item Bypass Network
\end{itemize}

\subsection*{Execution Interface}
Dispatches instructions to the appropriate execution unit based on the decoded operation.
Supported execution units include:
\begin{itemize}
  \item ALU
  \item Branch Unit
  \item Load/Store Unit
  \item Multiply/Divide Unit
  \item SIMD Unit
  \item Floating Point Unit
\end{itemize}

\subsection*{Commit Interface}
Transfers completed instructions to the Graduation List and Commit stage where architectural
state is updated.

\subsection*{Recovery Interface}
Provides rollback support whenever
\begin{itemize}
  \item branch prediction fails,
  \item exceptions occur,
  \item interrupts arrive.
\end{itemize}
Recovery information is exchanged with the Rename stage and Commit logic to restore a
precise architectural state.

\section{Internal Architecture}

Internally, the datapath acts as the interconnection fabric of the processor.
Instead of implementing computation directly, it instantiates and connects all major pipeline
components.

The internal architecture consists of:
\begin{itemize}
  \item Pipeline Registers
  \item Instruction Queue
  \item Rename Tables
  \item Physical Register File
  \item Bypass Network
  \item Functional Units
  \item Commit Logic
\end{itemize}

\section{Pipeline Registers}

To achieve high throughput, Sargantana separates adjacent pipeline stages using dedicated
pipeline registers.

Important pipeline registers include:
\begin{itemize}
  \item F1-to-F2
  \item F2-to-D
  \item IR-to-RR
  \item RR-to-EXE
  \item EXE-to-WB
\end{itemize}

These registers isolate combinational logic between stages, allowing multiple instructions to
progress simultaneously through different portions of the processor.
